\documentclass[11pt,a4paper,twoside]{tesis}
% SI NO PENSAS IMPRIMIRLO EN FORMATO LIBRO PODES USAR
%\documentclass[11pt,a4paper]{tesis}

\usepackage{graphicx}
\usepackage[utf8]{inputenc}
\usepackage[T1]{fontenc}
\usepackage{lmodern}
\usepackage[spanish,es-tabla,es-nodecimaldot]{babel}
\usepackage[left=3cm,right=3cm,bottom=3.5cm,top=3.5cm]{geometry}
\usepackage{amsmath}
\usepackage{xcolor}
\usepackage{listings}

%%%% LISTADOS DE CÓDIGO (Solidity y Haskell)

\lstdefinelanguage{Solidity}{
  morekeywords={contract, mapping, function, public, private, internal, external,
    constructor, require, assert, revert, return, returns, if, else, uint, int,
    address, bool, string, payable, view, pure, memory, pragma, solidity, enum,
    struct, msg, true, false},
  sensitive=true,
  morecomment=[l]{//},
  morecomment=[s]{/*}{*/},
  morestring=[b]",
}

% Las anotaciones de Liquid Haskell ({-@ ... @-}) no se muestran como comentarios.
\lstdefinelanguage{LHaskell}[]{Haskell}{deletecomment=[n]{\{-}{-\}}}

\lstset{
  basicstyle=\ttfamily\small,
  keywordstyle=\bfseries,
  commentstyle=\itshape\color{gray},
  stringstyle=\color{black!70},
  columns=fullflexible,
  keepspaces=true,
  showstringspaces=false,
  breaklines=true,
  frame=lines,
  framesep=4pt,
  xleftmargin=1em,
  aboveskip=1em,
  belowskip=1em,
  captionpos=b,
}
\renewcommand{\lstlistingname}{Listado}

\begin{document}

%%%% CARATULA

\def\autor{Tomás Jarabroviski}
\def\tituloTesis{Liquid Types para la verificación de \vspace{.2cm} \\ \emph{smart contracts} con diccionarios}
\def\runtitulo{Liquid Types para la verificación de smart contracts con diccionarios}
\def\runtitle{Liquid Types for the verification of smart contracts with dictionaries}
\def\director{Dr. Hernán Melgratti}
% Sin codirector: no definir \codirector.
\def\lugar{Buenos Aires, 2026}
\newcommand{\HRule}{\rule{\linewidth}{0.2mm}}
%
\thispagestyle{empty}

\begin{center}\leavevmode

\vspace{-2cm}

\begin{tabular}{l}
\includegraphics[width=2.6cm]{logofcen.pdf}
\end{tabular}


{\large \sc Universidad de Buenos Aires

Facultad de Ciencias Exactas y Naturales

Departamento de Computaci\'on}

\vspace{6.0cm}

%\vspace{3.0cm}
%{
%\Large \color{red}
%\begin{tabular}{|p{2cm}cp{2cm}|}
%\hline
%& Pre-Final Version: \today &\\
%\hline
%\end{tabular}
%}
%\vspace{2.5cm}

\begin{huge}
\textbf{\tituloTesis}
\end{huge}

\vspace{2cm}

{\large Tesis de Licenciatura en Ciencias de la Computaci\'on}

\vspace{2cm}

{\Large \autor}

\end{center}

\vfill

{\large

{Director: \director}

\vspace{.2cm}

% El codirector es opcional: solo se muestra si se define \codirector.
\ifdefined\codirector
{Codirector: \codirector}

\vspace{.2cm}
\fi

\lugar
}

\newpage\thispagestyle{empty}


%%%% ABSTRACTS, AGRADECIMIENTOS Y DEDICATORIA
% TODO: todavía tienen el texto de ejemplo de la plantilla; se escriben al final.
\frontmatter
\pagestyle{empty}
%%\begin{center}
%\large \bf \runtitulo
%\end{center}
%\vspace{1cm}
\chapter*{\runtitulo}

\noindent La princesa Leia, líder del movimiento rebelde que desea reinstaurar la República en la galaxia en los tiempos ominosos del Imperio, es capturada por las malévolas Fuerzas Imperiales, capitaneadas por el implacable Darth Vader. El intrépido Luke Skywalker, ayudado por Han Solo, capitán de la nave espacial ``El Halcón Milenario'', y los androides, R2D2 y C3PO, serán los encargados de luchar contra el enemigo y rescatar a la princesa para volver a instaurar la justicia en el seno de la Galaxia (aprox. 200 palabras).

\bigskip

\noindent\textbf{Palabras claves:} Guerra, Rebelión, Wookie, Jedi, Fuerza, Imperio (no menos de 5).

%\cleardoublepage
%%\begin{center}
%\large \bf \runtitle
%\end{center}
%\vspace{1cm}
\chapter*{\runtitle}

\noindent In a galaxy far, far away, a psychopathic emperor and his most trusted servant -- a former Jedi Knight known as Darth Vader -- are ruling a universe with fear. They have built a horrifying weapon known as the Death Star, a giant battle station capable of annihilating a world in less than a second. When the Death Star's master plans are captured by the fledgling Rebel Alliance, Vader starts a pursuit of the ship carrying them. A young dissident Senator, Leia Organa, is aboard the ship \& puts the plans into a maintenance robot named R2-D2. Although she is captured, the Death Star plans cannot be found, as R2 \& his companion, a tall robot named C-3PO, have escaped to the desert world of Tatooine below. Through a series of mishaps, the robots end up in the hands of a farm boy named Luke Skywalker, who lives with his Uncle Owen \& Aunt Beru. Owen \& Beru are viciously murdered by the Empire's stormtroopers who are trying to recover the plans, and Luke \& the robots meet with former Jedi Knight Obi-Wan Kenobi to try to return the plans to Leia Organa's home, Alderaan. After contracting a pilot named Han Solo \& his Wookiee companion Chewbacca, they escape an Imperial blockade. But when they reach Alderaan's coordinates, they find it destroyed - by the Death Star. They soon find themselves caught in a tractor beam \& pulled into the Death Star. Although they rescue Leia Organa from the Death Star after a series of narrow escapes, Kenobi becomes one with the Force after being killed by his former pupil - Darth Vader. They reach the Alliance's base on Yavin's fourth moon, but the Imperials are in hot pursuit with the Death Star, and plan to annihilate the Rebel base. The Rebels must quickly find a way to eliminate the Death Star before it destroys them as it did Alderaan (aprox. 200 palabras).

\bigskip

\noindent\textbf{Keywords:} War, Rebellion, Wookie, Jedi, The Force, Empire (no menos de 5). % OPCIONAL: comentar si no se quiere

%\cleardoublepage
%\chapter*{Agradecimientos}

\noindent Lorem ipsum dolor sit amet, consectetur adipiscing elit. Fusce sapien ipsum, aliquet eget convallis at, adipiscing non odio. Donec porttitor tincidunt cursus. In tellus dui, varius sed scelerisque faucibus, sagittis non magna. Vestibulum ante ipsum primis in faucibus orci luctus et ultrices posuere cubilia Curae; Mauris et luctus justo. Class aptent taciti sociosqu ad litora torquent per conubia nostra, per inceptos himenaeos. Mauris sit amet purus massa, sed sodales justo. Mauris id mi sed orci porttitor dictum. Donec vitae mi non leo consectetur tempus vel et sapien. Curabitur enim quam, sollicitudin id iaculis id, congue euismod diam. Sed in eros nec urna lacinia porttitor ut vitae nulla. Ut mattis, erat et laoreet feugiat, lacus urna hendrerit nisi, at tincidunt dui justo at felis. Class aptent taciti sociosqu ad litora torquent per conubia nostra, per inceptos himenaeos. Ut iaculis euismod magna et consequat. Mauris eu augue in ipsum elementum dictum. Sed accumsan, velit vel vehicula dignissim, nibh tellus consequat metus, vel fringilla neque dolor in dolor. Aliquam ac justo ut lectus iaculis pharetra vitae sed turpis. Aliquam pulvinar lorem vel ipsum auctor et hendrerit nisl molestie. Donec id felis nec ante placerat vehicula. Sed lacus risus, aliquet vel facilisis eu, placerat vitae augue.
 % OPCIONAL: comentar si no se quiere

%\cleardoublepage
%\hfill \textit{A mi persona favorita.}
  % OPCIONAL: comentar si no se quiere

\cleardoublepage
\tableofcontents

\mainmatter
\pagestyle{headings}

%%%% ACA VA EL CONTENIDO DE LA TESIS

%%%% CAPÍTULO: SOPORTE DE MAPPINGS

\chapter{Soporte de mappings}
\label{cap:mappings}

Los \emph{mappings} (diccionarios) son una de las estructuras de datos centrales en
los contratos que manejan fondos: los saldos de cada cuenta en un token o los montos
depositados por cada usuario en un banco se guardan en un \emph{mapping} de
direcciones en enteros, \lstinline[language=Solidity]|mapping(address => uint)|.
En~\cite{Consoli25}, donde se presentó la traducción de Solidity a Haskell que este
trabajo extiende, la posibilidad de representar \emph{mappings} se estudió solo de
forma preliminar, mediante una exploración manual que no llegó a incorporarse a la
traducción automática. En este capítulo desarrollamos ese soporte de manera completa.

El problema tiene dos niveles. El primero es de especificación: definir en Haskell una
representación de los \emph{mappings} junto con tipos refinados que describan con
precisión el efecto de sus operaciones, es decir, qué valor se lee después de una
escritura, qué entradas no cambian y cómo evoluciona la suma de los valores. El segundo
es de traducción: hacer que el traductor exprese cada operación de Solidity sobre un
\emph{mapping} en términos de esa representación, de modo que Liquid Haskell pueda
verificar automáticamente las funciones del contrato que la usan. Los tipos refinados
resultan adecuados para ambos niveles, porque permiten expresar los invariantes de la
estructura y las garantías de cada operación como parte de los tipos, y delegar su
verificación a un \emph{solver} SMT~\cite{VSJ+14}.

A lo largo del capítulo utilizaremos como ejemplo el token mínimo del
Listado~\ref{lst:token}. El constructor asigna todo el suministro inicial a la cuenta
que despliega el contrato; \lstinline|transfer| mueve \lstinline|amount| unidades
desde el saldo de quien invoca la función (\lstinline|msg.sender|) hacia el de
\lstinline|to|, y \lstinline|burn| destruye \lstinline|amount| unidades del saldo de
quien la invoca. Ambas funciones exigen, mediante \lstinline|require|, que el saldo
del remitente sea suficiente.

\begin{lstlisting}[language=Solidity, caption={Token mínimo con un \emph{mapping} de saldos.}, label={lst:token}]
contract Token {
    mapping(address => uint) balances;

    constructor(uint initialSupply) {
        balances[msg.sender] = initialSupply;
    }

    function transfer(address to, uint amount) public {
        require(balances[msg.sender] >= amount);
        balances[msg.sender] -= amount;
        balances[to] += amount;
    }

    function burn(uint amount) public {
        require(balances[msg.sender] >= amount);
        balances[msg.sender] -= amount;
    }
}
\end{lstlisting}

\section{Semántica de un mapping}
\label{sec:mappings-semantica}

Un \emph{mapping} de Solidity se comporta como una función total de claves en
valores. La traducción debe preservar las siguientes propiedades:

\begin{description}
  \item[Valor por defecto.] Toda clave tiene un valor asociado: si nunca fue escrita,
    su valor es el valor por defecto del tipo, que para \lstinline|uint| es $0$. En el
    ejemplo, \lstinline|balances[to]| vale $0$ para cualquier cuenta que nunca recibió
    tokens. En consecuencia, un contrato no puede distinguir una clave que nunca fue
    escrita de una cuyo valor almacenado es $0$, y \lstinline|delete m[k]| equivale a
    \lstinline|m[k] = 0|.
  \item[Escrituras locales.] Las sentencias \lstinline|m[k] = e|, \lstinline|m[k] += e|
    y \lstinline|m[k] -= e| modifican únicamente la entrada \lstinline|k|; el resto
    del \emph{mapping} no cambia.
  \item[Rango de los valores.] Los valores de tipo \lstinline|uint| son enteros entre
    $0$ y $2^{256}-1$. La herramienta representa este rango con el predicado
    $\mathit{InRange}$, que tiene dos definiciones según el modo en que se ejecute:
    \begin{itemize}
      \item en el modo por defecto, $\mathit{InRange}(v) \equiv v \geq 0$, es decir,
        los enteros se interpretan como enteros matemáticos, sin cota superior;
      \item en el modo de detección de desbordamiento, que se presenta en un capítulo
        posterior, $\mathit{InRange}(v) \equiv 0 \leq v \leq 2^{256}-1$.
    \end{itemize}
    En este capítulo usamos el modo por defecto, de modo que la única restricción
    sobre un valor almacenado es que sea mayor o igual a $0$.
\end{description}

La última propiedad es la que convierte a los \emph{mappings} en una fuente de
errores. En \lstinline|burn|, el \lstinline|require| garantiza que el saldo resultante
no sea negativo; si se omitiera, cualquier cuenta podría destruir tokens que no tiene.
Queremos que la traducción transforme esta propiedad en una obligación de prueba,
de modo que Liquid Haskell rechace la versión sin \lstinline|require|.

Una alternativa natural sería representar el \emph{mapping} como una función de Haskell
\lstinline[language=LHaskell]|k -> Int|. Sin embargo, además de leer y escribir claves
individuales, nos interesa poder razonar sobre la suma de todos los valores (por
ejemplo, el total de tokens en circulación), lo cual requiere una representación
finita y explícita de las entradas.

\section{Especificación en Liquid Haskell}
\label{sec:mappings-representacion}

Las definiciones del Listado~\ref{lst:mapwithbalance} forman parte de una biblioteca
que se incluye en cada módulo generado que utiliza algún \emph{mapping}.

\begin{lstlisting}[language=LHaskell, caption={Representación de un \emph{mapping} (definiciones principales).}, label={lst:mapwithbalance}]
{-@ type Map k v = {m : [(k, v)] | isUniq m} @-}

{-@ reflect sumBids @-}
sumBids :: Map k Int -> Int
sumBids []            = 0
sumBids ((_, x) : xs) = x + sumBids xs

{-@ data MapWithBalance k = MB { mMap :: Map k {v:Int | InRange v}
                               , mBal :: {v:Int | v >= 0} } @-}

{-@ type BalancedMap k =
      {m : MapWithBalance k | sumBids (mMap m) == mBal m} @-}
\end{lstlisting}

Las definiciones se leen de la siguiente manera:

\begin{itemize}
  \item \lstinline[language=LHaskell]|Map k v| es una lista de pares clave--valor. El
    predicado \lstinline|isUniq| exige que ninguna clave aparezca dos veces, de modo que
    cada clave tiene a lo sumo un valor asociado.
  \item \lstinline[language=LHaskell]|sumBids| calcula la suma de los valores del mapa.
    La anotación \lstinline|reflect| permite usar esta función dentro de los tipos
    refinados.
  \item \lstinline[language=LHaskell]|MapWithBalance k| acompaña al mapa con un campo
    \lstinline|mBal| que guarda la suma de sus valores. Los valores del mapa deben
    cumplir el predicado $\mathit{InRange}$ definido en la
    Sección~\ref{sec:mappings-semantica}.
  \item \lstinline[language=LHaskell]|BalancedMap k| es el tipo que efectivamente usa la
    traducción: un \lstinline|MapWithBalance| cuyo campo \lstinline|mBal| coincide con
    la suma real de sus valores. Este invariante es parte del tipo, por lo que toda
    función que devuelva un \lstinline|BalancedMap| debe preservarlo, y Liquid Haskell
    lo verifica.
\end{itemize}

La lista de pares no aparece en el código generado para las funciones del contrato:
este solo accede al mapa a través de las tres operaciones de la
Tabla~\ref{tab:map-interfaz}.

\begin{table}[ht]
\begin{lstlisting}[language=LHaskell, frame=lines, aboveskip=0pt, belowskip=0pt]
empWB          :: {m:BalancedMap k | mBal m == 0}
currentBalance :: m:BalancedMap k -> key:k -> {v:Int | InRange v}
setWB          :: m:BalancedMap k -> key:k -> val:{Int | InRange val}
               -> {n:BalancedMap k |
                     currentBalance n key == val &&
                     mBal n == mBal m - currentBalance m key + val}
\end{lstlisting}
\caption{Interfaz de \lstinline|BalancedMap| (tipos simplificados).}
\label{tab:map-interfaz}
\end{table}

\begin{itemize}
  \item \lstinline|empWB| es el mapa vacío, cuya suma es $0$.
  \item \lstinline|currentBalance m key| devuelve el valor asociado a \lstinline|key|,
    o $0$ si la clave no está en el mapa. Es la operación que implementa la semántica
    de valor por defecto: internamente se distingue una clave ausente de una con valor
    $0$, pero \lstinline|currentBalance| devuelve lo mismo en ambos casos, tal como
    ocurre en Solidity.
  \item \lstinline|setWB m key val| asocia \lstinline|val| a \lstinline|key|, ya sea
    agregando la clave o reemplazando su valor anterior, y actualiza la suma: se resta
    el valor previo y se suma el nuevo. Su argumento \lstinline|val| debe cumplir
    $\mathit{InRange}$; esta precondición es la que permite detectar saldos negativos.
\end{itemize}

Los tipos de la Tabla~\ref{tab:map-interfaz} resumen las especificaciones de la
biblioteca, que están expresadas en términos de la pertenencia de la clave al mapa y
de la función \lstinline|get|. Para enunciar con precisión las garantías que se
obtienen de ellas usamos la siguiente notación: $m[k]$ denota
\lstinline|currentBalance m k|, $m[k \mapsto v]$ denota \lstinline|setWB m k v|, y
$\mathrm{dom}(m)$ es el conjunto de claves presentes en $m$. Entonces, para todo
mapa, todo par de claves y todo valor en rango,
\[
  \forall\, m : \mathit{BalancedMap}\ K,\quad
  \forall\, k, k' \in K,\quad
  \forall\, v \text{ tal que } \mathit{InRange}(v),
\]
se cumplen las siguientes propiedades:

\begin{description}
  \item[Lectura por defecto.] $k \notin \mathrm{dom}(m) \;\Rightarrow\; m[k] = 0$.\\
    Una clave ausente se lee como $0$; en particular, toda lectura sobre
    \lstinline|empWB| da $0$.
  \item[Lectura tras escritura.] $m[k \mapsto v][k] = v$.\\
    Después de escribir $v$ en la clave $k$, leer $k$ devuelve $v$.
  \item[Marco.] $k' \neq k \;\Rightarrow\; m[k \mapsto v][k'] = m[k']$.\\
    Escribir una clave no modifica el valor de las demás.
  \item[Suma.] $\mathit{mBal}(m[k \mapsto v]) = \mathit{mBal}(m) - m[k] + v$.\\
    La suma se actualiza restando el valor anterior de la clave y sumando el nuevo.
  \item[Invariante.] $\mathit{mBal}(m) = \sum_{(j,\,x) \in \mathit{mMap}(m)} x$.\\
    El campo \lstinline|mBal| coincide siempre con la suma de los valores del mapa;
    es el refinamiento que define al tipo \lstinline|BalancedMap|.
  \item[Cota.] $m[k] \leq \mathit{mBal}(m)$.\\
    Ningún valor supera la suma, ya que todos los valores son no negativos.
\end{description}

Estas garantías se establecen dentro de la propia biblioteca, mediante lemas
auxiliares sobre listas y conjuntos que no es necesario conocer para seguir la
traducción.

\section{Traducción de las operaciones}
\label{sec:mappings-traduccion}

Recordemos brevemente el esquema de la traducción base~\cite{Consoli25}. El estado
del contrato se representa con un registro de Haskell, y cada función de Solidity se
traduce en una función que recibe el estado actual (\lstinline|st|) y sus parámetros
y devuelve el nuevo estado. La función \lstinline|revert| tiene tipo
\lstinline[language=LHaskell]!{v:String | False} -> a!, por lo que Liquid Haskell solo
acepta un \lstinline|revert| si puede probar que es inalcanzable. Sobre ese esquema:

\begin{itemize}
  \item \textbf{Secuencia.} Una secuencia de sentencias se traduce como una cadena de
    estados intermedios \lstinline|st1|, \lstinline|st2|, \ldots{} ligados con
    \lstinline|let|; el último de ellos es el resultado de la función.
  \item \textbf{\lstinline|require|.} La sentencia \lstinline|require(c)| se traduce de
    dos formas: como precondición en el tipo refinado de la función y como la
    sentencia \lstinline[language=LHaskell]|if c then st else revert "error"|, que
    continúa si la condición se cumple y revierte si no. La precondición es la que
    permite probar que ese \lstinline|revert| es inalcanzable.
  \item \textbf{\lstinline|assert|.} La sentencia \lstinline|assert(c)| se traduce
    como \lstinline[language=LHaskell]|if c then st else revert "assert"|, en el lugar
    donde está escrita. Como en este caso no hay una precondición que lo respalde,
    Liquid Haskell debe probar que \lstinline|c| se cumple siempre en ese punto de la
    función. Además, los \lstinline|assert| con los que termina la función se agregan
    como postcondiciones de su tipo refinado. Esto modifica la traducción
    de~\cite{Consoli25}, donde todo \lstinline|assert| se traducía únicamente como
    postcondición, de modo que un \lstinline|assert| escrito en medio de la función se
    verificaba sobre el estado final y no sobre el estado en el que está escrito.
\end{itemize}

Con estas reglas como base, la Tabla~\ref{tab:map-traduccion} muestra cómo se traducen las
operaciones sobre un \emph{mapping} \lstinline|m| declarado en un contrato
\lstinline|C|. Al igual que el resto de los campos del estado, su nombre lleva como
sufijo el nombre del contrato.

\begin{table}[ht]
\centering
\small
\begin{tabular}{ll}
\hline
\textbf{Solidity} & \textbf{Haskell} \\
\hline
\lstinline[language=Solidity]|mapping(K => uint) m| & campo \lstinline[language=LHaskell]|mC :: BalancedMap K| \\
valor inicial & \lstinline[language=LHaskell]|empWB| \\
\lstinline[language=Solidity]|m[k]| & \lstinline[language=LHaskell]|currentBalance (mC st) k| \\
\lstinline[language=Solidity]|m[k] = e| & \lstinline[language=LHaskell]|setWB (mC st) k e| \\
\lstinline[language=Solidity]|m[k] += e| & \lstinline[language=LHaskell]|setWB (mC st) k (currentBalance (mC st) k + e)| \\
\lstinline[language=Solidity]|m[k] -= e| & \lstinline[language=LHaskell]|setWB (mC st) k (currentBalance (mC st) k - e)| \\
\hline
\end{tabular}
\caption{Traducción de las operaciones sobre un \emph{mapping}. En las tres escrituras, el
  mapa resultante se asigna al campo \lstinline|mC| del nuevo estado.}
\label{tab:map-traduccion}
\end{table}

A continuación aplicamos la traducción al token del Listado~\ref{lst:token}.

\paragraph{Estado.} El estado del token se representa con el registro
\lstinline|TokenState| del Listado~\ref{lst:tokenstate}.

\begin{lstlisting}[language=LHaskell, caption={Estado del token.}, label={lst:tokenstate}]
data TokenState = TokenState {
  msgToken      :: Message,
  blockToken    :: Block,
  balanceToken  :: Amount,
  addressToken  :: Address,
  balancesToken :: BalancedMap Address
}
\end{lstlisting}

Los cuatro primeros campos se agregan en todo contrato: el mensaje de la invocación
en curso (que contiene \lstinline|msg.sender| y \lstinline|msg.value|), los datos del
bloque, la cantidad de ether que posee el contrato y su propia dirección. El último
campo es el \emph{mapping} \lstinline|balances|.

\paragraph{Constructor.} El constructor construye el estado inicial completo a partir
del mensaje y el bloque del despliegue. El \emph{mapping} parte del mapa vacío y, como
el constructor escribe una clave, se le aplica la escritura correspondiente:

\begin{lstlisting}[language=LHaskell]
{-@ constructorToken :: Message -> Block
                     -> {initialSupply:Uint | True}
                     -> {v:TokenState | True} @-}
constructorToken msgctx block initialSupply = TokenState
  { msgToken      = msgctx
  , blockToken    = block
  , balanceToken  = 0
  , addressToken  = fromString "Token"
  , balancesToken = setWB empWB (sender msgctx) initialSupply }
\end{lstlisting}

\noindent A diferencia del resto de las funciones, el constructor no recibe un estado
previo: recibe el mensaje (\lstinline|msgctx|) y el bloque (\lstinline|block|) del
despliegue, y su tipo refinado no tiene precondiciones porque no contiene
\lstinline|require|. El contrato comienza sin ether, porque el constructor no es
\lstinline|payable|, y su dirección se obtiene a partir de su nombre. La obligación que
impone \lstinline|setWB| es que \lstinline|initialSupply| sea no negativo, lo cual se
cumple porque el parámetro tiene tipo \lstinline|Uint|.

\paragraph{\lstinline|burn|.} El Listado~\ref{lst:burn} muestra el código generado
para \lstinline|burn|. Para facilitar la lectura abreviamos \lstinline|balancesToken|
como \lstinline|bal|, el remitente \lstinline|sender (msgToken st)| como
\lstinline|me| (es el mismo en todos los estados intermedios) y el mapa actualizado
\lstinline|balancesTokenNext| como \lstinline|next|; por lo demás, el código es el que
produce la herramienta.

\begin{lstlisting}[language=LHaskell, caption={Traducción de \lstinline|burn|.}, label={lst:burn}]
{-@ burn :: {st:TokenState | True}
         -> {amount:Uint | currentBalance (bal st) me >= amount}
         -> {v:TokenState | True} @-}
burn st amount = let
  st1 = if currentBalance (bal st) me >= amount
          then st else revert "error"
  st2 = let next = setWB (bal st1) me (currentBalance (bal st1) me - amount)
        in assert (lemma_currentBalance_get (bal st1) me) $
           assert (lemma_currentBalance_get next me) $
           st1 { bal = next }
  in st2
\end{lstlisting}

El \lstinline|require| aparece dos veces: como precondición y como la sentencia que
define \lstinline|st1|. La precondición menciona al parámetro \lstinline|amount|, por
lo que se escribe como refinamiento de ese parámetro y no del estado: en general, una
condición que menciona parámetros se asocia al último de ellos en la firma, de modo
que todos estén ligados al leerla. La postcondición es \lstinline|True|, porque la
función no contiene \lstinline|assert|.

La escritura \lstinline|balances[msg.sender] -= amount| se traduce en
\lstinline|st2|. Para aceptar la llamada a \lstinline|setWB|, Liquid Haskell debe
probar que el nuevo valor, \lstinline|currentBalance (bal st1) me - amount|, es no
negativo. Esto se deduce de la precondición, ya que \lstinline|st1| es igual a
\lstinline|st|. Las dos aserciones que acompañan a la escritura se explican al final
de la sección.

Si se elimina el \lstinline|require| de \lstinline|burn|, la precondición desaparece y
esa obligación ya no puede probarse: Liquid Haskell rechaza la llamada a
\lstinline|setWB| y la herramienta reporta el contrato como inseguro. El error de
programación (permitir que una cuenta destruya tokens que no tiene) se manifiesta
como un error de tipos en el código traducido.

\paragraph{\lstinline|transfer|.} Con las mismas abreviaturas, el tipo refinado de
\lstinline|transfer| es

\begin{lstlisting}[language=LHaskell]
{-@ transfer :: {st:TokenState | True}
             -> {to:Address | True}
             -> {amount:Uint | currentBalance (bal st) me >= amount}
             -> {v:TokenState | True} @-}
\end{lstlisting}

\noindent y tiene la misma forma que el de \lstinline|burn|: la condición del
\lstinline|require| queda asociada a \lstinline|amount|, el único parámetro que
menciona, y la postcondición es \lstinline|True| porque la función tampoco contiene
\lstinline|assert|. Las dos escrituras se traducen en dos estados consecutivos
(omitimos las aserciones de lemas, que son las mismas que en \lstinline|burn|):

\begin{lstlisting}[language=LHaskell]
  st2 = st1 { bal = setWB (bal st1) me (currentBalance (bal st1) me - amount) }
  st3 = st2 { bal = setWB (bal st2) to (currentBalance (bal st2) to + amount) }
\end{lstlisting}

La primera escritura genera la misma obligación que en \lstinline|burn|, que se
cumple gracias al \lstinline|require|. La segunda obligación es inmediata: la suma de
dos valores no negativos es no negativa. Notemos que la segunda escritura lee el
saldo de \lstinline|to| a partir de \lstinline|st2|, es decir, después de la primera
escritura. De esta forma la traducción es correcta incluso cuando \lstinline|to| y
\lstinline|msg.sender| son la misma cuenta, sin necesidad de tratar ese caso por
separado.

Supongamos ahora que el programador se equivoca de clave y escribe
\lstinline|require(balances[to] >= amount)|. La precondición pasa a hablar del saldo de
\lstinline|to|, que no dice nada sobre el saldo de \lstinline|msg.sender|; la
obligación de la primera escritura deja de poder probarse, y la herramienta reporta
nuevamente el contrato como inseguro. En cambio, la versión original del
Listado~\ref{lst:token} es verificada por Liquid Haskell.

\paragraph{Lemas en el código generado.} En el Listado~\ref{lst:burn}, la escritura
está acompañada de dos aserciones de la forma
\lstinline[language=LHaskell]|assert (lemma ...) $ ...|. La función
\lstinline[language=LHaskell]|assert :: a -> b -> b| ignora su primer argumento y
devuelve el segundo; su único propósito es incorporar al contexto de Liquid Haskell el
enunciado del lema. En este caso el lema es \lstinline|lemma_currentBalance_get m k|,
que relaciona \lstinline|currentBalance| con la función \lstinline|get|: si
$k \in \mathrm{dom}(m)$, entonces $m[k]$ es el valor que devuelve \lstinline|get|, y
si no, $m[k] = 0$. Como las especificaciones de \lstinline|setWB| están expresadas en
términos de \lstinline|get|, este lema es el que permite trasladarlas a las lecturas
con \lstinline|currentBalance| que aparecen en el código del contrato.

El traductor inserta los lemas según una regla fija. En cada escritura
\lstinline|m[k] = e|, \lstinline|m[k] += e| o \lstinline|m[k] -= e|:

\begin{itemize}
  \item siempre, \lstinline|lemma_currentBalance_get| sobre el mapa anterior y sobre el
    mapa resultante, para la clave \lstinline|k|;
  \item en el modo de detección de desbordamiento, además, para cada otra clave del mismo \emph{mapping} que lea la función, el lema de
    marco \newline\lstinline|lemma_currentBalance_setWB_maybe_other|;
  \item si el \emph{mapping} registra el ether recibido por el contrato, además, el
    lema de cota \lstinline|lemma_currentBalance_le_mBal| para la clave escrita y para
    cada clave del mismo \emph{mapping} que lea la función.
\end{itemize}

\noindent Las dos últimas situaciones se presentan en capítulos posteriores. En el
token, que se verifica en el modo por defecto y no recibe ether, solo aparecen las dos
aserciones del primer caso.

% ---------------------------------------------------------------------------
% PENDIENTE (párrafo de cierre quitado a pedido; recuperar más adelante):
%  - Alcance: el soporte cubre mappings con valores uint declarados como
%    variables de estado. Quedan fuera los mappings con valores de otros tipos
%    (bool, estructuras), los mappings usados como parámetros o variables
%    locales, los mappings anidados (mapping(address => mapping(address => uint)),
%    usados p. ej. para los permisos de gasto de un token) y la sentencia delete;
%    las sentencias con mappings anidados o delete se reportan como no traducidas.
%    -> sección de alcance/limitaciones.
%  - Puente a solvencia: el campo mBal y la garantía de cota, que acá solo
%    sostienen el invariante del mapa, permiten relacionar los saldos de un
%    mapping con el ether que posee el contrato. -> capítulo de solvencia.
% ---------------------------------------------------------------------------


%%%% BIBLIOGRAFIA
\backmatter
% TODO: quitar \nocite{*} cuando todas las referencias estén citadas en el texto.
\nocite{*}
\bibliographystyle{alpha}
\bibliography{tesis}

\end{document}
